\documentclass[a4paper,10pt]{report}

% \def\theme{nuit}

\newcommand\NomAppli        {Table}
\newcommand\NomAppliAccent  {complète}
\newcommand\TypeDocument    {Cahier des charges}
\newcommand\SousTitre       {Application Android de suggestion de repas, d'accords mets-boissons et de liste de courses.}
\newcommand\VersionDoc      {0.1}
\newcommand\DateDoc         {\today}
\newcommand\Plateforme      {Android · Jetpack Compose}
\newcommand\SourcesDonnees  {TheMealDB / TheCocktailDB}
\newcommand\Etablissement   {ESGI · M1}
\newcommand\Auteurs         {Mamadou Cisse · Theo Robert}
\newcommand\Encadrant       {M. Ramzy Kermad}

\input{style/01-paquets}
\input{style/02-couleurs}
\input{style/03-typographie}
\input{style/04-composants}
\input{style/05-mise-en-page}
\input{style/06-titres}
\input{style/07-listes-tableaux-figures}
\input{style/08-exigences}
\input{style/09-page-de-garde}
\input{style/10-liens}


\begin{document}

\pageDeGarde
\tableofcontents

\chapter{Spécimen du style}

\section{Typographie}

Texte courant en Onest : l'application propose chaque jour une carte
composée d'un plat, d'un verre et d'un dessert, en tenant compte de la
météo, du régime et du matériel de l'utilisateur. Chiffres :
350~g · 25~min · 4~pers. · 12,5~\%~vol. \textbf{Gras}, \emph{emphase},
\textmoyen{Onest Medium}, \textsecondaire{texte secondaire},
\texttt{MealRepository.kt}, \enquote{citation} et un
\href{https://www.themealdb.com}{lien}.

\subsection{Sous-section}

\subsubsection{Sous-sous-section}

\paragraph{Paragraphe} Texte qui suit directement le titre de paragraphe.

\section{Signatures}

\bandeau{Carte du jeudi 5 novembre}{13° éclaircies}

\medskip
\etiquette[1]{Le plat}\quad
\badge{25 min · sans four}\quad
\badge[herbe]{Adapté · sans gluten}\quad
\badge[sauge]{Sans alcool}\quad
\badge[contour]{Plaques}\quad
\badge[encre]{Japon}

\begin{lignes}
  \ligne[2]{Le verre}{Vermentino}
  \ligne[3]{Le dessert}{Affogato}
  \ligne{Farfalle}{350 g}
\end{lignes}

\begin{carte}[Le bandeau daté]
  En-tête de carte : la date et la météo cadrent la suggestion du jour.
  Filet épais dessus, filet fin dessous.
\end{carte}

\begin{note}
  Les données proviennent de TheMealDB et TheCocktailDB.
\end{note}

\begin{attention}[Allergènes]
  Indicatif : vérifiez toujours les étiquettes.
\end{attention}

\section{Exigences}

\begin{exigence}[priorite=haute]{Afficher la carte du jour}
  \label{ex:carte}
  L'écran d'accueil présente un plat, un verre et un dessert, accompagnés
  du bandeau daté (date et météo locale).
\end{exigence}

\begin{exigence}[type=ENF, priorite=moyenne]{Mode sombre}
  L'application suit le thème du système (palettes Blanc / Nuit).
\end{exigence}

\begin{exigence}[priorite=basse]{Carte du monde}
  Parcourir les recettes par pays d'origine.
\end{exigence}

Renvoi vers l'exigence~\ref{ex:carte}.

\section{Listes et tableaux}

\begin{itemize}
  \item Carte du jour
  \item Explorer : plats, cocktails, vins
  \begin{itemize}
    \item filtres par pays
  \end{itemize}
\end{itemize}

\begin{enumerate}
  \item Vérification d'âge
  \item Onboarding : régime, matériel, boissons
  \item Accueil
\end{enumerate}

\begin{description}
  \item[Accord] association d'un plat et d'une boisson.
\end{description}

\begin{table}[htbp]
  \caption{Écrans principaux}
  \begin{tableau}{L{3.2cm} Y R{2cm}}
    \entete{Écran} & \entete{Rôle} & \entete{Priorité} \\
    \midrule
    Carte & Suggestion du jour & \badge[herbe]{Haute} \\
    \midrule
    Explorer & Catalogue plats, cocktails, vins & \badge[agrume]{Moyenne} \\
    \midrule
    Courses & Liste générée depuis les recettes & \badge[herbe]{Haute} \\
  \end{tableau}
\end{table}

\ecran{1 · Accueil, la carte du jour.png}{Accueil — la carte du jour (Direction A, à remplacer par la C).}

\end{document}
